\section{Introduction}
\label{sec:introduction}

Scientific publications contain rich information beyond conventional
bibliographic metadata. For experimental studies, this includes information
about the measurements and experimental conditions underlying the reported
results. Much of this information, however, remains embedded in publications
rather than available as structured metadata. Recovering it could support
scientific discovery, reuse, reproducibility, and research data management.

This experimental information is often relational. A paper may describe
multiple samples, measurement datasets, and facilities, each with its own
attributes and links to other entities. Recovering such a record therefore
requires more than extracting isolated values: a system must identify the
entities reported in the paper, distinguish them at the appropriate scientific
granularity, and recover the links among them.

Photon and neutron science provides a useful setting for studying this
metadata extraction problem. FAIR data management and the capture and
interlinking of experimental metadata are established priorities across the
PaN community
\citep{murphy2025fair,expands2022fair,lohstroh2025daphne}.

Recent work has demonstrated the potential of schema-guided LLMs for
extracting metadata from scientific publications \citep{alyafeai2025mole},
while structured-extraction benchmarks have introduced field-sensitive metrics
and semantic alignment for evaluating heterogeneous structured outputs
\citep{ferguson2026extractbench}. However, publication-derived experimental
records containing repeated scientific entities and cross-entity references
remain comparatively underexplored.

To study this problem, we introduce a 40-paper benchmark for
publication-derived PaN metadata extraction. A two-layer schema combines
generic publication metadata with domain-specific experimental metadata,
representing facilities, samples, and measurement datasets as repeated
entities, with dataset records referencing their associated facilities and
samples. Ten papers formed the calibration set used to refine the schema,
field instructions, and annotation protocol. The remaining 30 formed the
evaluation set used for pairwise annotation agreement and all reported model 
comparisons.

We develop a benchmark-specific entity-aware evaluation methodology. It 
aligns predicted and gold entities, separates entity recovery from attribute 
quality among aligned entities, and remaps output-local identifiers before 
scoring cross-entity references.

We evaluate seven LLMs on the 30 evaluation papers. In the full-instruction 
zero-shot baseline, author and affiliation recovery is consistently high, 
whereas recovery of facilities, samples, and datasets is lower and model-dependent. 
In the instruction-ablation comparison, detailed field instructions yield higher 
sample and dataset recovery across all seven models. Patterns for model size and 
few-shot demonstrations are less uniform. Higher entity recovery does not necessarily 
correspond to higher conditional attribute or reference quality.
